
\documentclass[template=tabling,81pt,headonall]{azmoon}
\usepackage{xepersian}
\usepackage{amsfonts}
\usepackage{graphicx}
\usepackage{svg}
\svgpath{ {./images/} }
\graphicspath{ {./images/} }
\settextfont{Yas}
\setdigitfont{A Iranian Sans}
\usepackage{fontawesome5}

\printanswers
    \teacher{محمد صالح علی اکبری}
    \teachertitle{دبیر}
    \city{گناباد}
    \schooltitle{هنرستان}
    \school{خوارزمی}
    \grade{دهم}
    \branch{۱}
    \topic{ریاضی}
    \examdate{شهریور ۱۴۰۵}
    \answertime{۶۰ دقیقه}
    \begin{document}
	\begin{questions}
		\nointerlineskip%
		\vskip-\baselineskip
		\section{پودمان ۱: نسبت و تناسب}\question[1]{%
اگر $4$ کارگر بتوانند ساخت یک دیوار کارگاهی را در مدت $6$ روز به پایان برسانند، $8$ کارگر با همان توان کاری، این دیوار را در چند روز می‌سازند؟
    \begin{fourchoice}[4]\choice{۳ روز}
\choice{۱۲ روز}
\choice{۸ روز}
\choice{۲ روز}
\end{fourchoice}

            }\question[8]{%
نسبت طول به عرض یک صفحۀ فلزی در کارگاه برابر $\dfrac{5}{3}$ است. اگر عرض این صفحه $12\text{ cm}$ باشد:
    \begin{parts}[1]\part{طول این صفحه را به دست آورید.}
\part{محیط و مساحت این صفحه را محاسبه کنید.}
\end{parts}
‌
\\‌
\\
    }\question[8]{%
در سیستم اندازه‌گیری انگلیسی، هر اینچ معادل $2/54$ سانتی‌متر است.
    \begin{parts}[1]\part{طول یک قطعۀ صنعتی برابر $10$ اینچ است؛ این اندازه معادل چند سانتی‌متر است؟}
\part{طول یک تختۀ چوب $127$ سانتی‌متر است؛ طول آن را بر حسب اینچ محاسبه کنید.}
\end{parts}
‌
\\‌
\\
    }\section{پودمان ۲: درصد و کاربردهای آن}\question[2]{%
قیمت اولیۀ یک ابزار کارگاهی ۱۰۰ هزار تومان بوده است. اگر قیمت آن ابتدا $\%20$ افزایش و سپس $\%20$ کاهش یابد، قیمت نهایی چند هزار تومان خواهد بود؟
    \begin{fourchoice}[4]\choice{۱۰۰ هزار تومان}
\choice{۹۶ هزار تومان}
\choice{۱۰۴ هزار تومان}
\choice{۹۸ هزار تومان}
\end{fourchoice}

            }\question[8]{%
قیمت اصلی یک لباس کار برابر ۴۰۰ هزار تومان است. اگر فروشگاهی $\%15$ تخفیف برای آن در نظر بگیرد:
    \begin{parts}[1]\part{مبلغ تخفیف چقدر است؟}
\part{خریدار برای خرید این لباس چه مبلغی باید پرداخت کند؟}
\end{parts}
‌
\\‌
\\
    }\question[8]{%
تولید روزانۀ یک کارگاه از $200$ قطعه به $250$ قطعه افزایش یافته است.
    \begin{parts}[1]\part{مقدار افزایش تولید چند قطعه بوده است؟}
\part{درصد تغییر (درصد رشد) تولید را به دست آورید.}
\end{parts}
‌
\\‌
\\
    }\section{پودمان ۳: معادله‌های درجۀ دوم}\question[2]{%
جواب‌های معادلۀ درجۀ دوم $x^2 - 5x + 6 = 0$ کدام گزینه‌ها هستند؟
    \begin{fourchoice}[4]\choice{$x = 2$ و $x = 3$}
\choice{$x = -2$ و $x = -3$}
\choice{$x = 1$ و $x = 6$}
\choice{$x = -1$ و $x = -6$}
\end{fourchoice}

            }\question[8]{%
معادلۀ درجۀ دوم $x^2 - 4x - 5 = 0$ را با استفاده از روش دلتا ($\Delta$) حل کرده و ریشه‌های آن را بیابید.‌
\\}\question[8]{%
طول یک باغچۀ مستطیل‌شکل $3$ متر بیشتر از عرض آن است. اگر مساحت باغچه برابر $28\text{ m}^2$ باشد:
    \begin{parts}[1]\part{با در نظر گرفتن عرض باغچه برابر $x$، یک معادلۀ درجۀ دوم برای مسئله بنویسید.}
\part{عرض و طول باغچه را با حل معادله بیابید.}
\end{parts}
‌
\\
    }\section{پودمان ۴: توان‌رسانی به توان عددهای گویا}\question[2]{%
حاصل عبارت عددی $8^{\frac{1}{3}}$ کدام است؟
    \begin{fourchoice}[4]\choice{$2$}
\choice{$4$}
\choice{$16$}
\choice{$64$}
\end{fourchoice}

            }\question[8]{%
حاصل هر یک از عبارت‌های زیر را به ساده‌ترین صورت به دست آورید:
    \begin{parts}[2]\part{$\sqrt[3]{27} + \sqrt[2]{16} = $}
\part{$\sqrt[3]{-125} = $}
\end{parts}
‌
\\
    }\question[8]{%
عبارت تواندار زیر را با استفاده از قوانین پایه‌های مساوی ساده کنید و پاسخ را به صورت یک توان کسری بنویسید: $A = x^{\frac{3}{2}} \times x^{\frac{1}{2}} \quad (x > 0)$‌
\\}\section{پودمان ۵: نسبت‌های مثلثاتی}\question[1]{%
در یک مثلث قائم‌الزاویه، طول ضلع مقابل به زاویۀ $\theta$ برابر ۳ و طول وتر برابر ۵ است. مقادیر $\sin(\theta)$ و $\tan(\theta)$ به ترتیب کدامند؟
    \begin{fourchoice}[4]\choice{$\dfrac{3}{5}$ و $\dfrac{3}{4}$}
\choice{$\dfrac{4}{5}$ و $\dfrac{4}{3}$}
\choice{$\dfrac{3}{5}$ و $\dfrac{4}{3}$}
\choice{$\dfrac{4}{5}$ و $\dfrac{3}{4}$}
\end{fourchoice}

            }\question[1.5]{%
در یک مثلث قائم‌الزاویه، طول وتر $10\text{ cm}$ و یکی از زاویه‌های حاده $30^\circ$ است.
    \begin{parts}[1]\part{طول ضلع مقابل به این زاویه را بیابید ($\sin 30^\circ = \dfrac{1}{2}$).}
\part{طول ضلع مجاور به این زاویه را بیابید ($\cos 30^\circ = \dfrac{\sqrt{3}}{2}$).}
\end{parts}

    }\question[1.5]{%
یک نردبان به طول $6$ متر به دیوار کارگاهی تکیه داده شده است و با سطح زمین زاویۀ $60^\circ$ می‌سازد.
    \begin{parts}[1]\part{ارتفاع نقطه‌ای از دیوار که نردبان به آن تکیه کرده چقدر است؟ ($\sin 60^\circ = \dfrac{\sqrt{3}}{2}$).}
\part{فاصلۀ پای نردبان تا دیوار چقدر است؟ ($\cos 60^\circ = \dfrac{1}{2}$).}
\end{parts}

    }\end{questions}
    \end{document}
    